\exercisesection{希尔伯特–塞缪尔级数}

\begin{enumerate}
\def\labelenumi{\arabic{enumi})}
\tightlist
\item
  令 \(\Gamma\) 为交换群，H 为 \(\Gamma\)-型分次环（A, II，第~164页）。
\end{enumerate}

对于 \(\gamma\in\Gamma\)，记 H \(_{\gamma}\) 为 H 的次数为 \(\gamma\) 的齐次分量。假设 H 由 H \(_{0}\) 和有限族 (x \(_{i}\) ) \(_{i\in I}\) 的次数 \(\neq0\) 的齐次元素生成。记 \(\gamma_{i}\) 为 x \(_{i}\) 的次数。对于每个交换群 K，记 K{[}{[}Γ{]}{]} 为乘积群 K \(^{\Gamma}\)，并将 \(\sum_{\gamma\in\Gamma}m_{\gamma}\text{T}^{\gamma}\) 表示为元素 (m \(_{\gamma}\) )，它属于 K \(^{\Gamma}\)。记 K{[}Γ{]} 为 K \(^{(\Gamma)}\) 在 K{[}{[}Γ{]} 中的子群。赋予 K{[}{[}Γ{]}{]} 以其作为 Γ 的代数 Z{[}Γ{]}-模的自然结构。最后，令 C 为 H₀-模类的一个遗传集合（A, VIII, § 10, 第 1 号, 定义 1），令 K(C) 为相应的 Grothendieck 群（同上，第 2 号）。对于每个 C 型 H₀-模 E，记 {[}E{]} 为其在 K(C) 中的类。

\begin{enumerate}
\def\labelenumi{\alph{enumi})}
\tightlist
\item
  假设 \(\mathbf{H}_0\) 是诺特的，并令 \(\mathbf{M}\) 为有限生成的分次 H-模，使得对于每个 \(\gamma \in \Gamma\)，\(\mathbf{H}_0\)-模 \(\mathbf{M}_{\gamma}\) 属于 C 型。证明元素 \(\prod_{i \in I} (1 - T^{\gamma_i}) \sum_{\gamma \in \Gamma} [\mathbf{M}_{\gamma}] T^{\gamma}\) 属于
\end{enumerate}

\(K(\mathbb{C})[[\Gamma]]\) 属于 \(K(\mathbb{C})[\Gamma]\)。

\begin{enumerate}
\def\labelenumi{\alph{enumi})}
\setcounter{enumi}{1}
\tightlist
\item
  现在假设 \(\Gamma = \mathbf{Z}^s\)，\(\mathrm{H}_{\gamma} = 0\)，且 \(\mathbf{M}_{\gamma} = 0\) 对 \(\gamma \notin \mathbf{N}^s\) 成立，并且 \(\gamma_i\) 是 \(\mathbf{Z}^s\) 的基向量；记 \(a_k\) 为这样的 \(\gamma_i\) 的数目：它们等于第 \(k\) 个基向量，证明有
\end{enumerate}

\[
\prod_ {k = 1} ^ {s} (1 - \mathrm{T} _ {k}) ^ {a _ {k}} \sum_ {\gamma \in \mathbf {Z} ^ {s}} [ \mathrm{M} _ {\gamma} ] \mathrm{T} ^ {\gamma} = \sum_ {\gamma} m _ {\gamma} \mathrm{T} ^ {\gamma}
\]

其中 \(\mathbf{T} = (\mathrm{T}_{1},\dots,\mathrm{T}_{s})\)，且 \((m_{\gamma})\) 是有限支撑族。推出存在元素 \(c_{\boldsymbol{b}}\in\mathrm{K}(\mathbb{C})\)，其中 \(\boldsymbol{b}\in\mathbf{Z}^{s}\)，除有限多个外都为零，使得当 \(\gamma\in\mathbf{N}^{s}\) 充分大时，有

\[
[ \mathbf {M} _ {\gamma} ] = \sum_ {\boldsymbol {b}} c _ {\boldsymbol {b}} \mathbf {B} _ {\boldsymbol {b}} (\gamma)
\]

其中

\[
\mathrm{B} _ {b} (\gamma) = \prod_ {j = 1} ^ {s} \frac {(\gamma_ {j} + 1) \dots (\gamma_ {j} + b _ {j})}{b _ {j} !},
\]

求和遍历 \(Z^{s}\) 中满足 \(b_{i} < a_{i}\)（对 \(1 \leqslant i \leqslant s\)）的元素 b。

\begin{enumerate}
\def\labelenumi{\arabic{enumi})}
\setcounter{enumi}{1}
\tightlist
\item
  令 \(f: \mathbf{Z} \to \mathbf{Z}\) 为映射。证明以下三个性质等价：
\end{enumerate}

\begin{enumerate}
\def\labelenumi{\alph{enumi})}
\item
  存在 \(n_0, n_1\) 属于 \(\mathbf{Z}\) 以及 \(\mathrm{P}\) 属于 \(\mathbf{Q}[\mathrm{T}]\)，使得 \(f(n) = 0\) 当 \(n \leqslant n_1\)，且 \(f(n) = \mathrm{P}(n)\) 当 \(n \geqslant n_0\)。
\item
  存在一个有限支撑族 \((a_i)_{i\in \mathbf{Z}}\)，其元素属于 \(\mathbf{Z}\)，使得 \(f(n) = \sum_{i\in \mathbf{Z}}a_i\left[ \begin{array}{cc}n + i\\ i \end{array} \right]\) 对每个 \(n\in \mathbf{Z}\) 都成立。
\item
  存在 \(r \in \mathbf{Z}\) 和 \(Q \in \mathbf{Z}[T, T^{-1}]\)，使得 \(\sum_{n \in \mathbf{Z}} f(n) T^n = (1 - T)^r Q(T)\)。
\end{enumerate}

于是称 \(f\) 在无穷远处为多项式。

推广到从 Z 到任意 Z-模的映射情形。

\begin{enumerate}
\def\labelenumi{\arabic{enumi})}
\setcounter{enumi}{2}
\tightlist
\item
  \begin{enumerate}
  \def\labelenumii{\alph{enumii})}
  \tightlist
  \item
    举出一个域 K 上的 Z-型分次代数 H 的例子：它由有限多个次数 \textgreater{} 0 的齐次元素生成，且映射 \(n \mapsto \dim_{\mathbf{K}} \mathrm{H}_{n}\) 在无穷远处不是多项式。
  \end{enumerate}
\end{enumerate}

\begin{enumerate}
\def\labelenumi{\alph{enumi})}
\setcounter{enumi}{1}
\tightlist
\item
  采用第 4 节第 2 号定理 1 的记号和假设，令 D 为 \(d_{i}\) 的最小公倍数。证明 H 是环 \(\mathrm{H}_{0}[(\mathrm{X}_{i}^{\mathrm{D}/d_{i}})_{i\in\mathbb{I}}]\) 上的有限生成模。
\end{enumerate}

推出对于每个 \(n_{0}\in\mathbf{N}\)，级数 \(\sum_{n\in\mathbf{Z}}\mathrm{long}_{\mathrm{H}_{0}}(\mathbf{M}_{n_{0}+n\mathbf{D}})\mathbf{T}^{n}\) 具有形式 \(\frac{\mathbf{P}_{n_{0}}(\mathbf{T})}{(1-\mathbf{T})^{c}}\)，其中 \(c=\mathrm{Card}(\mathbf{I})\) 且 \(\mathbf{P}_{n_{0}}(\mathbf{T})\in\mathbf{Q}[\mathbf{T},\mathbf{T}^{-1}]\)。推出映射 \(n\mapsto\mathrm{long}_{\mathrm{H}_{0}}(\mathbf{M}_{n_{0}+n\mathbf{D}})\) 在无穷远处为多项式。

\begin{enumerate}
\def\labelenumi{\arabic{enumi})}
\setcounter{enumi}{3}
\tightlist
\item
  置身于第 4 节第 2 号定理 1 的情形，记 C 为分次 H-模的同构类集合，其中各分量作为 \(\mathbf{H}_0\)-模具有有限长度，并记相应的 Grothendieck 群为 \(\mathrm{K}(\mathbb{C})\)。
\end{enumerate}

\begin{enumerate}
\def\labelenumi{\alph{enumi})}
\tightlist
\item
  对于每个上述 H-模 M，证明存在有限长度的 \(H_{0}\)-分次模 \(M_{0}\) 以及满射分次同态 \(p:H\otimes_{H_{0}}M_{0}\to M\)。由此推出每个上述 H-模 M 都有一个由 H-模组成的分次分解（A, X，第~56页）：
\end{enumerate}

\[
\dots \rightarrow \mathrm{H} \otimes_ {\mathrm{H} _ {0}} \mathrm{M} _ {1} \rightarrow \mathrm{H} \otimes_ {\mathrm{H} _ {0}} \mathrm{M} _ {0} \rightarrow \mathrm{M} \rightarrow 0
\]

其中对所有 \(i\)，\(M_{i}\) 都是有限长度的 \(H_{0}\)-分次模。

\begin{enumerate}
\def\labelenumi{\alph{enumi})}
\setcounter{enumi}{1}
\item
  使用 H-模 \(\mathbf{H}_0\) 的 Koszul 分解（A, X，第~152页），证明对于每个分次 H-模 M，\(\mathrm{Tor}_j^{\mathrm{H}}(\mathrm{H}_0,\mathrm{M})\) 都自然带有一个分次；且 \(\mathrm{Tor}_j^{\mathrm{H}}(\mathrm{H}_0,\mathrm{M}) = 0\) 对 \(j > \mathrm{Card}(\mathrm{I})\) 成立；并且 \(\mathrm{Tor}_j^{\mathrm{H}}(\mathrm{H}_0,\mathrm{P}\otimes_{\mathrm{H}_0}\mathrm{H})\) 是具有有限长度的 \(\mathbf{H}_0\)-模（对每个 \(j\)），其中 P 是具有有限长度的 \(\mathbf{H}_0\)-模。
\item
  证明对于每个 \(j\)，当 M 为 C 型时，\(\mathrm{Tor}_{j}^{\mathrm{H}}(\mathrm{H}_{0}, \mathrm{M})\) 是具有有限长度的 \(\mathrm{H}_{0}\)-分次模。（可对 \(j\) 作归纳来证明。若 \(j = 0\)，使用 a) 和 b)。一般情形下，使用正合序列 \(0 \to \mathbf{N} \to \mathbf{H} \otimes_{\mathrm{H}_{0}} \mathbf{M}_{0} \to \mathbf{M} \to 0\)、归纳假设和 b)。）
\item
  对于每个 C 型 H-模 M，置
\end{enumerate}

\[
e (\mathbf {M}) = \sum_ {p \in \mathbf {Z}} \sum_ {j = 0} ^ {\infty} (- 1) ^ {j} [ \operatorname{Tor} _ {j} ^ {\mathrm{H}} (\mathrm{H} _ {0}, \mathbf {M}) _ {p} ] \mathrm{T} ^ {p},
\]

  从而得到 \(\mathbf{K}(\mathcal{C}_{0})[\mathrm{T},\mathrm{T}^{-1}]\) 中的一个元素，其中 \(\mathbf{K}(\mathcal{C}_{0})\) 是有限长度 \(\mathrm{H}_{0}\)-模的 Grothendieck 群。证明 \(\mathbf{M}\mapsto e(\mathbf{M})\) 在分次 H-模的正合序列上可加，从而 \(\mathbf{M}\mapsto e(\mathbf{M})\) 确定一个群同态

\[
e: \mathrm{K} (\mathfrak {C}) \to \mathrm{K} (\mathfrak {C} _ {0}) [ \mathrm{T}, \mathrm{T} ^ {- 1} ].
\]

对于每个有限长度的 \(H_{0}\)-分次模 P，置

\[
\tau (\mathrm{P}) = [ \mathrm{H} \otimes_ {\mathrm{H} _ {0}} \mathrm{P} ] \in \mathrm{K} (\mathbb {C}).
\]

证明 \(\mathbf{P} \mapsto \tau(\mathbf{P})\) 在正合序列上可加，从而 \(\mathbf{P} \mapsto \tau(\mathbf{P})\) 确定相应 Grothendieck 群之间的一个同态

\[
\tau : \mathrm{K} (\mathbb {C} _ {0}) [ \mathrm{T}, \mathrm{T} ^ {- 1} ] \rightarrow \mathrm{K} (\mathbb {C}).
\]

证明 \(e\) 与 \(\tau\) 互为逆群同构。（首先证明 \(e \circ \tau\) 是恒等映射。为证明这一断言，只需证明 \(\tau\) 是满射；也就是说，只需证明对于每个 C 型 H-模 M，{[}M{]} 属于 \(\tau\) 的像。注意到 M 有一个有限滤过，其相邻商被 \(H_0\) 的某个极大理想湮没；并且一个被 \(H_0\) 的极大理想 m 湮没的 C 型 H-模 M 有一个由有限生成自由 H/mH-模组成的有限分次分解（A, X，第~56页）。）

\begin{enumerate}
\def\labelenumi{\alph{enumi})}
\setcounter{enumi}{4}
\tightlist
\item
  令 M 为 C 型 H-模。证明在 K(C) 中有
\end{enumerate}

\[
[ \mathbf {M} ] = \sum_ {i \geqslant 0} (- 1) ^ {i} [ \mathrm{H} \otimes_ {\mathrm{H} _ {0}} \operatorname{Tor} _ {i} ^ {\mathrm{H}} (\mathrm{H} _ {0}, \mathbf {M}) ].
\]

\begin{enumerate}
\def\labelenumi{\alph{enumi})}
\setcounter{enumi}{5}
\tightlist
\item
  令 \(\Gamma\) 为交换群，\(\Phi : K(\mathbb{C}) \to \Gamma\) 为群同态。对于每个有限长度的 \(H_0\) -模 P，置
\end{enumerate}

\[
\psi (\mathrm{P}) = \Phi ([ \mathrm{H} \otimes_ {\mathrm{H} _ {0}} \mathrm{P} ]).
\]

证明对于每个 C 型 H-模 M，有

\[
\Phi ([ \mathbf {M} ]) = \sum_ {j \geqslant 0} (- 1) ^ {j} \psi (\operatorname{Tor} _ {j} ^ {\mathrm{H}} (\mathrm{H} _ {0}, \mathbf {M})).
\]

特别地，采用第 4 节第 2 号备注 4 的记号，证明有

\[
\mathrm{Q} _ {\mathbf {M}} = \sum_ {j = 0} ^ {r} (- 1) ^ {j} \mathrm{P} _ {\mathrm{T} _ {j}},
\]

并从而有

\[
c _ {\mathbf {M}} = \sum_ {j = 0} ^ {r} (- 1) ^ {j} \mathrm{long} _ {\mathrm{H} _ {0}} (\mathrm{Tor} _ {j} ^ {\mathrm{H}} (\mathrm{H} _ {0}, \mathbf {M})).
\]

\begin{enumerate}
\def\labelenumi{\alph{enumi})}
\setcounter{enumi}{6}
\tightlist
\item
  使用前面的结果，给出第 4 节第 1 号定理 1 的一个新证明。
\item
  令 Γ 为交换群。用从 Specmax(H₀) × Z 到 Γ 的映射描述所有定义在 C 型 H-模上、取值于 Γ 的可加映射 M ↦ Φ(M)。
\item
  假设 H₀ 是域。令 A 为环。描述所有可加映射 \(M\mapsto C(M)\)，其取值于 \(A\)，并满足 \(C(M)\cdot C(N) = \sum_j(-1)^j C(\operatorname{Tor}_{j}^{\mathrm{H}}(M,N))\)，其中 M 和 N 为任意一对有限生成分次 H-模。
\end{enumerate}

\begin{enumerate}
\def\labelenumi{\arabic{enumi})}
\setcounter{enumi}{4}
\tightlist
\item
  \begin{enumerate}
  \def\labelenumii{\alph{enumii})}
  \tightlist
  \item
    令 H 为 N-型分次环，满足 \(H_{0}\) 是域且 H 是有限生成的 \(H_{0}\)-代数。令 M 和 N 为两个有限生成分次 H-模。置 \(T_{i} = \operatorname{Tor}_{i}^{\mathrm{H}}(\mathbf{M}, \mathbf{N})\)。证明等式
  \end{enumerate}
\end{enumerate}

\[
\sum_ {i = 0} ^ {\infty} (- 1) ^ {i} \mathrm{P} _ {\mathrm{T} _ {i}} = \frac {\mathrm{P} _ {\mathrm{M}} \cdot \mathrm{P} _ {\mathrm{N}}}{\mathrm{P} _ {\mathrm{H}}}.
\]

\begin{enumerate}
\def\labelenumi{\alph{enumi})}
\setcounter{enumi}{1}
\tightlist
\item
  令 A 为诺特局部环，M 和 N 为两个 A-模。证明：若对于 \(m \in Z\) 和 \(i \geqslant 1\)，\(\operatorname{Tor}_{i}^{\operatorname{gr}_{m}(A)}(\operatorname{gr}_{m}(M), \operatorname{gr}_{m}(N))\) 为零，则 \(\operatorname{Tor}_{i}^{A}(M, N) = 0\) 对于 \(i \geqslant 1\) 成立，并且 \(H_{M \otimes_{A} N, m} = \frac{H_{M, m} \cdot H_{N, m}}{H_{A, m}}\) 对于 \(m \in Z\) 成立。
\end{enumerate}

\begin{enumerate}
\def\labelenumi{\arabic{enumi})}
\setcounter{enumi}{5}
\tightlist
\item
  令 H 为 Z-型分次环，满足 \(H_{n} = \{0\}\) 对 n \textless{} 0 成立，且 \(H_{0}\) 是 Artin 局部环，H 是有限生成的 \(H_{0}\)-代数。令 L. 为有限生成自由分次 H-模的有界复形（参见 A, X，第~56页）。于是对每个 i，\(L_{i}\) 同构于形如 \(\bigoplus_{j=1}^{r_{i}} H(-n_{ij})\) 的 H-模，其中 \(n_{ij} \in Z\)，而 \(H(-k)\) 是满足 \(H(-k)_{n} = H_{n-k}\) 的分次 H-模。对于每个有限生成分次 H-模 M，记 \(G_{M} \in Q[T]\) 为唯一多项式，使得存在整数 \(n_{1} \geqslant 0\) 满足
\end{enumerate}

\[
\mathrm{G} _ {\mathrm{M}} (n) = \sum_ {i <   n} \operatorname{long} _ {\mathrm{H} _ {0}} (\mathrm{M} _ {i}), \quad \text {for} \quad n \geqslant n _ {1}.
\]

证明关系式

\[
\sum_ {i \geqslant 0} (- 1) ^ {i} \mathrm{G} _ {\mathrm{H} _ {i} (\mathrm{L.})} = \sum_ {k \geqslant 0} \frac {(- 1) ^ {k}}{k !} \rho_ {k} \frac {\partial^ {k} \mathrm{G} _ {\mathrm{H}}}{\partial \mathrm{T} ^ {k}}
\]

其中 \(\rho_{k} = \sum_{i\geqslant 0}(-1)^{i}\sum_{j = 1}^{r_{i}}n_{ij}\)，\(n_{ij}\) 是与 \(\mathbf{L}_i\) 相关的次数，且 \(\mathrm{H}_i(\mathrm{L.})\) 表示复形 \(\mathrm{L.}\) 的同调。

¶ 7) a) 令 l 为整数 ≥ 1，且 n ∈ N。证明存在唯一的递减整数序列：\(a_{l,l}(n) \geqslant a_{l,l-1}(n) \geqslant \cdots \geqslant a_{l,1}(n) \geqslant 0\)，使得

\[
n = \sum_ {i = 1} ^ {l} \binom {a _ {l, i} (n) + i - 1} {i}.
\]

证明 \(n \leqslant m\)，当且仅当按照字典序，\((a_{l,l}(n), a_{l,l-1}(n), ..., a_{l,1}(n))\) 小于或等于 \((a_{l,l}(m), ..., a_{l,1}(m))\)。

置

\[
\partial_ {l} (n) = \sum_ {i = 1} ^ {l} \binom {a _ {l, i} (n) + i} {i + 1},
\]

并且当 \(l \geqslant 2\) 时，记 \(\partial_l^*(n)\) 为最小整数 \(p\)，使得 \(\partial_{l-1}(p) \geqslant n\)。

求 \(a_{l+1,i}(\partial_l(n)), a_{l-1,i}(\partial_l^*(n))\) 关于 \(a_{l,i}(n)\) 的表达式。

\begin{enumerate}
\def\labelenumi{\alph{enumi})}
\setcounter{enumi}{1}
\tightlist
\item
  若映射 H:N → N 满足
\end{enumerate}

\[
\left\{ \begin{array}{l l} \mathrm{H} (0) = 1, \\ \mathrm{H} (l + 1) \leqslant \partial_ {l} (\mathrm{H} (l)) & \text {for all} \quad l \geqslant 1. \end{array} \right.
\]

证明这等价于

\[
\left\{ \begin{array}{l} \mathrm{H} (0) = 1, \\ \partial_ {l} ^ {*} \mathrm{H} (l) \leqslant \mathrm{H} (l - 1) \quad \text {for all} \quad l \geqslant 2. \end{array} \right.
\]

令 H : N → N 为 Macaulay 函数。证明只可能有两种情形：α) 存在 r ∈ N，使得对所有 j ≥ r 都有 H(j) = 0。

β) 存在 \(r \in \mathbf{N}\) 和一个有限的整数序列

\[
b _ {0} \geqslant b _ {1} \geqslant \dots \geqslant b _ {r} \geqslant 0
\]

使得对于所有 \(j \geqslant r\)，都有

\[
(*) \quad \mathrm{H} (j) = \sum_ {n = 0} ^ {r} \binom {b _ {n} - n + j} {b _ {n}}.
\]

证明 \(r\) 和序列 \(b_0 \geqslant b_1 \geqslant \cdots \geqslant b_r \geqslant 0\) 由函数 \(H\) 唯一确定。置 \(d(H) = 0\) 于情形 \(\alpha\)，置 \(d(H) = b_0 + 1\) 于情形 \(\beta\)，于是 \(d(H) - 1\) 是关于变量 \(j\) 的多项式 (*) 的次数。令 \(a(H) \frac{j^{d(H)-1}}{(d(H) - 1)!}\) 为多项式 (*) 的最高次项。证明 \(a(H)\) 是整数 \(> 0\)。

证明：一个关于变量 \(j\)、系数属于 \(\mathbf{Q}\) 的多项式，若在整数上取整数值且对足够大的 \(j\) 严格为正，则不一定具有形式 (*)。

\begin{enumerate}
\def\labelenumi{\alph{enumi})}
\setcounter{enumi}{2}
\tightlist
\item
  令 H: N → N 为 Macaulay 函数。置
\end{enumerate}

\[
\mathrm{S} _ {\mathrm{H}} (\mathrm{T}) = \sum_ {n = 0} ^ {\infty} \mathrm{H} (n) \mathrm{T} ^ {n},
\]

从而 \(S_{H}(T) \in Z[[T]]\)。证明 \(S_{H}(T) \in Z[T]\)，或者存在一个有限的整数序列

\[
c _ {0} \geqslant c _ {1} \geqslant \dots \geqslant c _ {r} > 0
\]

该序列由 H 唯一确定，并存在多项式 \(P \in Z[T]\)，使得

\[
(* *) \mathrm{S} _ {\mathrm{H}} (\mathrm{T}) = \sum_ {n = 0} ^ {r} \frac {\mathrm{T} ^ {n}}{(1 - \mathrm{T}) ^ {c _ {n}}} + \mathrm{P}.
\]

证明 \(c_0 = d(\mathbf{H})\)，且 \(a(\mathbf{H})\) 是整数 \(c_i\) 的个数，其中 \(c_i = d(\mathbf{H})\)。

有理分式 \(\sum_{n=0}^{r}\frac{T^n}{(1-T)^{c_n}}\) 称为 \(H\) 的渐近展开。令 \(A_1, ..., A_d\) 为整数，满足 \(A_d \neq 0\)，使得有理分式

\[
\sum_ {n = 1} ^ {d} \frac {\mathrm{A} _ {n}}{(1 - \mathrm{T}) ^ {n}} - \mathrm{S} _ {\mathrm{H}} (\mathrm{T})
\]

是多项式。证明序列 \((\mathrm{A}_{n})_{1\leq n\leq d}\) 由 \(H\) 的渐近展开唯一确定，并且唯一决定 H 的渐近展开。特别地，有 \(d=d(\mathrm{H})\)，\(\mathrm{A}_{d}=a(\mathrm{H})\)。更详细地考察 \(d(\mathrm{H})=2\) 的情形。

\begin{enumerate}
\def\labelenumi{\alph{enumi})}
\setcounter{enumi}{3}
\tightlist
\item
  令 \(\mathbf{H}:\mathbf{N}\to\mathbf{N}\) 为 Macaulay 函数。证明以下条件等价：
\end{enumerate}

\begin{enumerate}
\def\labelenumi{(\roman{enumi})}
\item
  \(\partial_l^*\mathrm{H}(l) = \mathrm{H}(l - 1)\) 对每个 \(l\geqslant 2\)
\item
  H 是具有与 H 相同渐近展开的最小 Macaulay 函数。
\end{enumerate}

这样的 Macaulay 函数称为极值的。令 \(c_{0} \geqslant c_{1} \geqslant \cdots \geqslant c_{r} > 0\) 为整数序列。证明存在唯一一个渐近展开为 \(\sum_{n=0}^{r} \frac{\mathrm{T}^{n}}{(1 - \mathrm{T})^{c_{n}}}\) 的极值 Macaulay 函数。证明若 H 是极值 Macaulay 函数，则 \(\mathrm{H}(1) = d(\mathrm{H})\) 或 \(\mathrm{H}(1) = d(\mathrm{H}) + 1\)。由此推出对于每个 Macaulay 函数 H，都有 \(\mathrm{H}(1) \geqslant d(\mathrm{H})\)。

\begin{enumerate}
\def\labelenumi{\alph{enumi})}
\setcounter{enumi}{4}
\tightlist
\item
  令 H: N → N 为 Macaulay 函数。证明以下性质等价：
\end{enumerate}

\begin{enumerate}
\def\labelenumi{(\roman{enumi})}
\tightlist
\item
  有 \(\mathrm{S_H(T)} = \frac{1}{(1 - T)^{d(H)}}\)
\end{enumerate}

\begin{enumerate}
\def\labelenumi{(\roman{enumi})}
\setcounter{enumi}{1}
\item
  \(\mathbf{H}(1) = d(\mathbf{H})\)
\item
  \(\partial_{l}\mathbf{H}(l) = \mathbf{H}(l + 1)\) 对每个 \(l\geqslant 1\)
\end{enumerate}

特别地，任何满足 \(\mathrm{H}(1)=d(\mathrm{H})\) 的函数都是极值的。

\begin{enumerate}
\def\labelenumi{\arabic{enumi})}
\setcounter{enumi}{7}
\tightlist
\item
  令 E 为集合。置 \(\mathcal{M}(\mathrm{E}) = \mathbf{N}^{(\mathrm{E})}\)，并将 \(\mathcal{M}(\mathrm{E})\) 中的运算记为乘法。一个 \(\mathcal{M}(\mathrm{E})\) 的理想是 \(I\) 的子集 \(\mathcal{M}(\mathrm{E})\)，满足 \(I\cdot\mathcal{M}(\mathrm{E}) \subset I\)。\(\mathcal{M}(\mathrm{E})\) 的阶梯是一个理想的补集。记 \(d: \mathcal{M}(\mathrm{E}) \to \mathbf{N}\) 为满足 \(d(x) = 1\) 对所有 \(x \in \mathrm{E}\) 都成立的唯一幺半群同态。令 \(\mathrm{A} \subset \mathcal{M}(\mathrm{E})\)。对每个 \(l \in \mathbf{N}\)，记 \(\mathrm{A}_l\) 为集合 \(\mathrm{A} \cap d^{-1}(l)\)。还记 \(\delta \mathrm{A}\) 为由 \(m \in \mathcal{M}(\mathrm{E})\) 组成的集合：对 E 中每个元素 \(x\)，若其整除 \(m\)，则 \(m / x \in \mathrm{A}\)；并记 \(d\mathrm{A}\) 为如下 \(m\) 的集合：存在 \(x \in \mathrm{E}\) 使得 \(mx \in \mathrm{A}\)。
\end{enumerate}

\begin{enumerate}
\def\labelenumi{\alph{enumi})}
\tightlist
\item
  证明对于 \(\mathcal{M}(E)\) 的每个子集 A，都有 \(d\delta A \subset A\) 且 \(A \subset \delta dA\)；若 A 和 B 是 \(\mathcal{M}(E)\) 的两个子集，则关系 \(dA \subset B\) 与 \(A \subset \delta B\) 等价。证明对于子集 \(\Delta\)（属于 \(\mathcal{M}(E)\)），以下三个性质等价：
\end{enumerate}

\begin{enumerate}
\def\labelenumi{(\roman{enumi})}
\item
  \(\Delta\) 是阶梯；
\item
  \(\Delta \supset d\Delta\)；
\item
  \(\delta \Delta \supset \Delta\)。
\end{enumerate}

\begin{enumerate}
\def\labelenumi{\alph{enumi})}
\setcounter{enumi}{1}
\tightlist
\item
  令 F 为 E 的子集，并将 \(\mathcal{M}(\mathbf{F})\) 识别为 \(\mathcal{M}(\mathbf{E})\) 中的典范像。证明 \(\mathcal{M}(\mathbf{F})\) 的子集是 \(\mathcal{M}(\mathbf{F})\) 的阶梯，当且仅当它是 \(\mathcal{M}(\mathbf{E})\) 的阶梯。证明对于阶梯 \(\Delta\)（在 \(\mathcal{M}(\mathbf{E})\) 中），以下条件等价：
\end{enumerate}

\begin{enumerate}
\def\labelenumi{(\roman{enumi})}
\item
  \(\Delta_l\) 对每个 \(l \in \mathbf{N}\) 都是有限集；
\item
  \(\Delta_{1}\) 是有限集；
\item
  存在 E 的有限子集 F，使得 \(\Delta\subset\mathcal{M}(F)\)。
\end{enumerate}

称这样的阶梯为有限型的。

\begin{enumerate}
\def\labelenumi{\alph{enumi})}
\setcounter{enumi}{2}
\tightlist
\item
  从现在起，假设 E = N，并置 \(\mathcal{M}(E) = \mathcal{M}\)。令 k 为域。记 R 为分次代数 \(\mathrm{K}^{(\mathcal{M})} = \mathrm{K}[(\mathrm{X}_i)_{i \in \mathbb{N}}]\)。将 M 识别为 \(\mathrm{K}[(\mathrm{X}_i)_{i \in \mathbb{N}}]\) 中的单项式集合；对于 M 的每个子集 A，记 \(k(A)\) 为由 A 生成的 R 的向量子空间。令 I 为 M 的理想，\(\Delta\) 为其补阶梯。则 \(k(I)\) 是 R 的理想，而 \(k(\Delta)\) 是 \(k(I)\) 的补空间。令 \(J \subset R\) 为 \(R\) 的分次理想。证明存在阶梯 \(\Delta \subset M\)，使得 \(k(\Delta)\) 是 \(J\) 的补空间。（在 \(M\) 上置逆字典序，规定关系 \(X_1^{a_1} \ldots X_n^{a_n} \prec X_1^{b_1} \ldots X_n^{b_n}\) 成立的条件是存在 \(0 \leqslant p \leqslant n\)，使得 \(a_j = b_j\) 对 \(j \textgreater{} p\) 成立，且 \(a_p < b_p\)。按归纳定义 \(m_1, \ldots, m_p, \ldots\)，它们属于 \(M\)，如下：对每个 \(p\)，\(m_p\) 是 \(M\) 中与 \(m_1, \ldots, m_{p-1}\) 模 \(J\) 线性无关的最小单项式。令 \(\Delta\) 为 \(m_i\) 的集合。证明 \(\Delta\) 是阶梯。）这样的阶梯是有限型的，当且仅当代数 \(R/J\) 在 \(k\) 上是有限型的。
\end{enumerate}

  沿用前一习题的记号。在 \(M\) 上置字典序。令 \(l \in \mathbf{N}\)。称 \(A\) 为 \(\mathcal{M}_{l}\) 的初始子集，若对每个 \(m \in \mathrm{A}\)，集合

\[
[ m ] = \{m ^ {\prime} \in \mathcal {M} _ {d (m)} | m ^ {\prime} <   m \}
\]

包含于 A。

\begin{enumerate}
\def\labelenumi{\alph{enumi})}
\item
  令 \(l \in \mathbf{N}\)（分别令 \(l \in \mathbf{N} - \{0\}\)）且令 \(\mathbf{A} \subset \mathcal{M}_l\) 为初始段。证明 \(\delta \mathbf{A}\)（分别为 \(d\mathbf{A}\)）是 \(\mathcal{M}_{l+1}\)（分别为 \(\mathcal{M}_{l-1}\)）的初始段。假设 \(\mathbf{A}\) 有最大元素 \(m = X_1^{a_1} \ldots X_n^{a_n}\)。确定 \(\delta \mathbf{A}\)（分别 \(d\mathbf{A}\)）的最大元素。
\item
  令 \(l \in \mathbf{N} - \{0\}\)，并令 \(\mathbf{A} \subset \mathcal{M}_l\) 为非空有限初始段。令 \(\mathrm{X}_{a_l}\) 为满足 \([\mathrm{X}_{a_l}^l] \subset \mathbf{A}\) 的最大变量。证明每个 \(m\)（属于 \(\mathbf{A} - [\mathrm{X}_{a_l}^l]\)）都等于 1 或被 \(\mathrm{X}_{a_l + 1}\) 整除，并且 \(\mathrm{A}' = \{m \in \mathcal{M}_{l-1}|m\mathrm{X}_{a_l + 1} \in \mathbf{A} - [\mathrm{X}_{a_l}^l]\}\) 是 \(\mathcal{M}_{l-1}\) 的初始段。由前述结果推出 \(\operatorname{Card}([X_1^{b_1} ... X_n^{b_n}])\) 的一个表达式。
\item
  由前述结果推出，对于每个 \(l \in \mathbf{N} - \{0\}\) 和 \(\mathcal{M}_l\) 的每个有限初始段 A，都有 \(\text{Card}(\delta\text{A}) = \partial_l(\text{Card}(\text{A}))\)，其中 \(\partial_l\) 在习题 7 中定义。类似地计算 \(\text{Card}(d\text{A})\)。
\item
  对于每个 \(l \in \mathbf{N}\) 和每个 \(n \in \mathbf{N}\)，置
\end{enumerate}

\[
\mu (n) = \inf_{\substack{\mathbf{A}\subset \mathcal{M}_{l}\\ \operatorname{Card}(\mathbf{A}) = n}}\operatorname{Card}(d\mathbf{A}),
\]

并记 \(m_{n}\) 为第 \(n\) 个单项式（按 \(\mathcal{M}_{l}\) 的字典序排列）（置 \(m_{0}=1\)）。

对于 \(\mathcal{M}_{l}\) 的每个有限子集 A，记 CA 为 \(\mathcal{M}_{l}\) 的初始段，使得 \(\mathrm{Card(A)} = \mathrm{Card(CA)}\)。证明以下三个断言等价：

(MC1) 对于每个 \(l \in \mathbf{N}\) 和每个子集 \(\mathbf{A} \subset \mathcal{M}_l\)，都有 \(\mathbf{C}\delta \mathbf{A} \subset \delta \mathbf{CA}\)；

(MC2) 对于每个 \(l \in \mathbf{N}\) 和每个子集 \(\mathbf{A} \subset \mathcal{M}_l\)，都有 \(d\mathbf{CA} \subset \mathbf{CdA}\)；

(MC3) 对于每个 \(l \in \mathbf{N}\) 和每个 \(n \in \mathbf{N}\)，都有 \(\operatorname{Card}(d[m_n]) = \mu(n)\)。

\begin{enumerate}
\def\labelenumi{\arabic{enumi})}
\setcounter{enumi}{9}
\tightlist
\item
  采用前两道习题的记号。令 \(l \in \mathbf{N}\) 且 \(\mathbf{A} \subset \mathcal{M}_l\)。对于每一对 \((i, d)\)，置
\end{enumerate}

\(\mathrm{A}_{(i;d)} = \{m \in \mathrm{A} | m \text{能被} \mathrm{X}_i^d \text{整除且不能被} \mathrm{X}_i^{d+1}\}\)

并记 \(C_{(i;d)}A\) 为 \(\operatorname{Card}(A_{(i;d)})\) 中按字典序排列的前 \((\mathcal{M}_{l})_{(i;d)}\) 个元素所成的集合（按字典序）。置 \(C_{i}A = \bigcup C_{(i;d)}A\)。

\begin{enumerate}
\def\labelenumi{\alph{enumi})}
\tightlist
\item
  对于每个 \(m \in \mathcal{M}_{l}\)，若存在，记 \(\pi(m)\) 为 \(\mathcal{M}_{l}\) 中 \(m\) 的逆字典序前驱（即 \(m \neq X_{1}^{l}\)）。令 \(i\) 为满足 \(1 < i\) 的整数，令 \(a_{i}\) 为严格正整数。证明有
\end{enumerate}

\[
\begin{array}{c} \pi (X _ {i} ^ {a _ {i}} \dots X _ {p} ^ {a _ {p}}) = X _ {i - 1} X _ {i} ^ {a _ {i} - 1} \dots X _ {p} ^ {a _ {p}} \\ \pi (X _ {1} ^ {a _ {1}} X _ {i} ^ {a _ {i}} \dots X _ {p} ^ {a _ {p}}) = X _ {i - 1} ^ {a _ {i - 1} + 1} X _ {i} ^ {a _ {i} - 1} \dots X _ {p} ^ {a _ {p}}. \end{array}
\]

\begin{enumerate}
\def\labelenumi{\alph{enumi})}
\setcounter{enumi}{1}
\item
  令 \(\mathbf{A} \subset \mathcal{M}_l\)，且 \(p \geqslant 4\)，使 \(\mathbf{A} \subset [\mathbf{X}_p^l]\)（即 \(\mathbf{A}\) 中的单项式只涉及变量 \(\mathbf{X}_j\)，\(1 \leqslant j \leqslant p\)）。假设对每个整数 \(i\) 满足 \(1 \leqslant i \leqslant p\)，都有 \(\mathbf{A} = \mathbf{C}_i\mathbf{A}\)。证明 \(\mathbf{A} = \mathbf{CA}\)（采用习题 9, d) 的记号）。
\item
  令 \(\mathbf{A} \subset \mathcal{M}_l\)，满足 \(\mathbf{A} \subset [\mathbf{X}_3^l]\)，且 \(\mathbf{A} = \mathbf{C}_i\mathbf{A}\) 对 \(i = 1, 2, 3\) 成立。置 \(\mathbf{A} = \bigcup_{k=0} \mathbf{B}_k\mathbf{X}_3^{k}\)，其中 \(\mathbf{B}_k \subset [\mathbf{X}_2^{l-k}]\)，且 \(b(k) = \text{Card}(\mathbf{B}_k)\)。证明 \(\mathbf{B}_k\) 是 \(\mathcal{M}_{l-k}\) 的初始段，并且 \(0 \leqslant b(0) \leqslant l + 1\)，有 \(b(k+1) < b(k)\)（当 \(b(k) \neq 0\) 时），且 \(k \mapsto b(k)\) 是递减的。考察该断言的逆命题。令 \(\mu\) 为整数 \(k\) 的个数，其中 \(b(k) = l + 1 - k\)。证明 \(\text{Card}(d\mathbf{A}) = (\text{Card}(\mathbf{A})) - \mu\)。
\end{enumerate}

¶ 11) 我们拟证明习题 9, d) 的断言 MC1、MC2、MC3（Macaulay 定理）。令 \(l \in \mathbf{N}\)、\(\mathbf{A} \subset \mathcal{M}_l\)，并令 \(p\) 为使 \(\mathbf{A} \subset [\mathbf{X}_p^l]\) 的最小整数。我们将对 \(p\) 作归纳来证明 MC2。

\begin{enumerate}
\def\labelenumi{\alph{enumi})}
\item
  证明 \(p = 1\) 和 \(p = 2\) 时定理成立。
\item
  利用归纳假设证明对于所有 \(i \leqslant p\)，都有 \(dC_{i}A \subset C_{i}dA\)（采用习题 8 的记号）。
\item
  置 \(\mathbf{A}^1 = \mathbf{A}\)，\(\mathbf{A}^{j+1} = \mathbf{C}_i\mathbf{A}^j\)，其中 \(i \equiv j \mod p\) 且 \(1 \leqslant i \leqslant p\)。证明当 \(j\) 充分大时有 \(\mathbf{A}^{j+1} = \mathbf{A}^j\)。（对于每个 \(a \in \mathcal{M}_l\)，令 \(n(a) \in \mathbf{N}\) 表示 \(a\) 在字典序中的秩，并置 \(n(\mathbf{A}) = \sum_{a \in \mathbf{A}} n(a)\)。证明 \(n(\mathbf{A}^{j+1}) \leqslant n(\mathbf{A}^j)\)，且 \(n(\mathbf{A}^{j+1}) = n(\mathbf{A}^j)\) 当且仅当 \(\mathbf{A}^{j+1} = \mathbf{A}^j\)。）
\item
  证明一般情形下的定理。（称 \(A \subset [X_{p}^{l}] \subset M_{l}\) 为极小的，如果 \(\operatorname{Card}(dA) \leqslant \operatorname{Card}(dA')\) 对满足 \(A' \subset [X_{p}^{l}]\) 且 \(\operatorname{Card}(A') = \operatorname{Card}(A)\) 的子集成立。只需证明若 A 是极小的，则 CA 也是极小的。令 A 为极小的。利用 b)，证明 \(C_{i}A\) 对所有 \(i \leqslant p\) 都是极小的。由 c) 的记号推出 \(A^{j}\) 对所有 j 都是极小的。由 c) 推出存在极小的 B，使得 \(\operatorname{Card}(B) = \operatorname{Card}(A)\)，且 \(C_{i}B = B\) 对 \(1 \leqslant i \leqslant p\) 成立。在 p = 3 的情形，使用习题 10, c) 得出结论。在 \(p \geqslant 4\) 的情形，使用习题 10, b) 推出 \(B = CA\)。
\end{enumerate}

¶ 12) 令 H : N → N 为映射，令 k 为域。我们拟证明以下性质的等价性：

\begin{enumerate}
\def\labelenumi{(\roman{enumi})}
\item
  H 是 Macaulay 函数（第~90页，习题 7）。
\item
  存在有限型阶梯（第~92页，习题 8）\(\Delta \subset \mathcal{M}\)，使得对所有 \(l \in \mathbf{N}\)，都有 \(\mathrm{H}(l) = \mathrm{Card}(\Delta_l)\)。
\item
  存在一个分次 \(k\)-代数 \(\mathbf{A}\)，其类型为 \(\mathbf{N}\)，由其有限多个次数为 1 的元素生成，使得对每个 \(l\)，都有 \(\mathrm{H}(l) = \dim_k A_l\)。
\end{enumerate}

\begin{enumerate}
\def\labelenumi{\alph{enumi})}
\item
  为证明 (ii) \(\Leftrightarrow\) (iii)，我们将使用习题 8, c)。
\item
  我们证明 (i) \(\Rightarrow\) (ii)。令 H 为 Macaulay 函数。对每个 \(l\)，令 \(\Delta_l \subset \mathcal{M}_l\) 为满足 \(\operatorname{Card}(\Delta_l) = \mathrm{H}(l)\) 的初始部分，并置 \(\Delta = \bigcup_{l} \Delta_l\)。为证明 \(\Delta\) 是阶梯，我们将使用习题 9, c)。
\item
  我们证明 (ii) \(\Rightarrow\) (i)。令 \(\Delta\) 为阶梯。为证明 \(\mathbf{H}: l \mapsto \operatorname{Card}(\Delta_l)\) 是 Macaulay 函数，我们将使用 Macaulay 定理（MC1）（习题 11）和习题 9, c) \(^{1}\)。
\end{enumerate}
% semantic-fragments: legacy-input-seam
\relax{}
